

\section{Preface}

In this book I set out to describe entity modelling from first principles.
Writing this book I have aimed to provide a careful, well-considered account. 
Like many books on this subject, I proceed by way of examples, but I have tried to avoid common pitfalls such as pseudo-definitions, where one term is explained only in terms of another that is no more clearly defined. Where possible, I have related the terminology, as it is introduced, to that of mathematics --- specifically to the terminologies used in set theory, logic, and category theory.

In the software development context, entity modelling is widely understood as a precursor to data definition. When used, it provides a firm conceptual base, but it is less often used than it should be because there has traditionally existed a discontinuity --- a gap --- between entity modelling and data definition. Traditional methodologies have advocated bridging this gap using a corrective step called data normalisation, but this approach --- modelling followed by normalisation --- reduces the residual value of the entity model, which is why entity modelling has remained a side activity rather than a key part of the software process.

This book aims to change that by presenting a version of entity modelling in which the transition to data definition is seamless.

Methodologically, this is achieved in part by paying attention to paths through the system of relationships within the subject domain and noting which comparable paths are actually equivalent.

In developing the exposition I also make use of terminology that has parallels in philosophical writing, in particular notions such as the universal and the particular. I do so not as a contribution to philosophy itself, but because these terms help to articulate distinctions that otherwise tend to be reinvented in different guises within software modelling. In this respect the aim is not terminological novelty but conceptual stability.

Following the approach described in this book eliminates the gap between entity modelling and data definition and renders data normalisation unnecessary but proving that this is so is an exercise in fairly abstract mathematics and is outside the scope of this book. It requires a level of mathematics that we do not wish to assume here.\footnote{A paper containing the proof is in preparation.}

It seems to me that the possibility of eliminating the entity modelling gap has not previously been recognised. As a result, software development has long borne the costs associated with either abandoning the documentation of conceptual structure or maintaining separate models and data definitions. This book sets out an account in which this cost is removed.


